% Part: first-order-logic
% Chapter: natural-deduction
% Section: soundness

\documentclass[../../../include/open-logic-section]{subfiles}

\begin{document}

\iftag{FOL}
      {\olfileid{fol}{ntd}{sou}}
      {\olfileid{pl}{ntd}{sou}}

\olsection{నిర్దుష్టత}

\begin{explain}
సహజ నిగమనం వంటి !!^a{derivation} వ్యవస్థ, వాస్తవంగా అనుగమించని
విషయాలను !!{derive} చేయలేకపోతే \emph{నిర్దుష్టమైనది}. అందువల్ల
నిర్దుష్టత అనేది !!{derivation} వ్యవస్థలకు ఒక రకమైన హామీగల భద్రతా
ధర్మం. పరిశీలిస్తున్న నిరూపణ-సిద్ధాంత ధర్మాన్ని బట్టి, ఉదాహరణకు
కిందివి తెలుసుకోవాలనుకుంటాం:
\begin{enumerate}
\item ప్రతి !!{derivable} !!{sentence} ఒక \iftag{FOL}{చెల్లుబాటు
  అయ్యేది}{సర్వసత్యం};
\item ఒక !!a{sentence} కొన్ని ఇతర వాక్యాల నుంచి !!{derivable} అయితే,
  అది వాటి పర్యవసానం కూడా;
\item ఒక !!{sentence}s సమితి వైరుధ్యభరితమైతే, అది అసంతృప్తమైనది.
\end{enumerate}
ఇవి ఒక !!a{derivation} వ్యవస్థకు ముఖ్యమైన ధర్మాలు. వీటిలో ఏదైనా
చెల్లకపోతే ఆ !!{derivation} వ్యవస్థ లోపభూయిష్టమైనది---అది మరీ ఎక్కువ
విషయాలను !!{derive} చేస్తుంది. కాబట్టి ఒక !!{derivation} వ్యవస్థ
నిర్దుష్టతను స్థాపించడం అత్యంత ముఖ్యం.
\end{explain}

\begin{thm}[నిర్దుష్టత]
\ollabel{thm:soundness}
!!{undischarged} పరికల్పనలు $\Gamma$ నుంచి $!A$ !!{derivable} అయితే,
$\Gamma \Entails !A$.
\end{thm}

\begin{proof}
$!A$ యొక్క ఒక !!a{derivation}ను $\delta$ అందాం. $\delta$లోని
నిగమనాల సంఖ్యపై ఆగమనం చేస్తాం.

ఆగమన ప్రాతిపదికలో నిగమనాల సంఖ్య $0$ అయినప్పుడు ప్రతిపాదనను చూపుతాం.
ఈ సందర్భంలో $\delta$ ఒకే !!{sentence}~$!A$, అంటే ఒక పరికల్పనను మాత్రమే
కలిగి ఉంటుంది. నిగమనాల ద్వారా మాత్రమే పరికల్పనలు !!{discharged}
కాగలవు; ఇక్కడ నిగమనాలే లేవు కనుక ఆ పరికల్పన !!{undischarged}గానే
ఉంటుంది. అందువల్ల నిరూపణలోని !!{undischarged} పరికల్పనలన్నిటినీ
సంతృప్తిపరచే ఏ \iftag{FOL}{!!{structure}~$\Struct{M}$}{!!{valuation}~$\pAssign{v}$}
అయినా $!A$ను కూడా సంతృప్తిపరుస్తుంది.

ఇప్పుడు ఆగమన మెట్టు. $\delta$లో $n$ నిగమనాలున్నాయని అనుకుందాం.
అన్నిటికంటే కింది నిగమనపు పూర్వాధారం(లు), ఒక్కొక్కదానిలో $n$ కంటే
తక్కువ నిగమనాలు గల ఉప-!!{derivation}sను ఉపయోగించి !!{derive} చేయబడ్డాయి.
ఆగమన పరికల్పనను అనుకుందాం: అన్నిటికంటే కింది నిగమనపు పూర్వాధారాలు,
ఆ పూర్వాధారాలతో ముగిసే ఉప-!!{derivation}sలోని !!{undischarged}
పరికల్పనల నుంచి అనుగమిస్తాయి. నిష్కర్ష~$!A$, మొత్తం నిరూపణలోని
!!{undischarged} పరికల్పనల నుంచి అనుగమిస్తుందని చూపాలి.

అన్నిటికంటే కింది నిగమనం రకాన్ని బట్టి సందర్భాలను వేరుచేస్తాం.
మొదట ఒకే పూర్వాధారం గల నిగమనాలను పరిశీలిస్తాం.

\begin{enumerate}
\item చివరి నిగమనం \Intro{\lnot} అనుకుందాం: !!{derivation} రూపం
  ఇలా ఉంటుంది:
  \begin{prooftree}
    \AxiomC{$\Gamma, \Discharge{!A}{n}$}
    \RightLabel{$\delta_1$}
    \DeduceC{$\lfalse$}
    \DischargeRule{\Intro{\lnot}}{n}
    \UnaryInfC{$\lnot !A$}
  \end{prooftree}
  ఆగమన పరికల్పన ప్రకారం $\delta_1$లోని !!{undischarged} పరికల్పనలు
  $\Gamma \cup \{!A\}$ నుంచి $\lfalse$ అనుగమిస్తుంది. ఒక
  \iftag{FOL}{!!a{structure}~$\Struct{M}$}{!!a{valuation}~$\pAssign{v}$}ను
  పరిశీలించండి. $\iftag{FOL}{\Sat{M}}{\pSat{v}}{\Gamma}$ అయితే
  $\iftag{FOL}{\Sat{M}}{\pSat{v}}{\lnot !A}$ అని చూపాలి. ప్రతిపక్ష
  నిరూపణ కోసం $\iftag{FOL}{\Sat{M}}{\pSat{v}}{\Gamma}$ కానీ
  $\iftag{FOL}{\Sat/{M}}{\pSat/{v}}{\lnot !A}$, అంటే
  $\iftag{FOL}{\Sat{M}}{\pSat{v}}{!A}$, అని అనుకుందాం. అప్పుడు
  $\iftag{FOL}{\Sat{M}}{\pSat{v}}{\Gamma \cup \{!A\}}$ అవుతుంది.
  ఇది మన ఆగమన పరికల్పనకు విరుద్ధం. కనుక
  $\iftag{FOL}{\Sat{M}}{\pSat{v}}{\lnot !A}$.

\item చివరి నిగమనం \Elim{\land}: దీనికి రెండు రూపాలున్నాయి:
  $!A \land !B$ అనే పూర్వాధారం నుంచి $!A$ లేదా $!B$ను నిగమించవచ్చు.
  మొదటి సందర్భాన్ని పరిశీలిద్దాం. !!{derivation}~$\delta$ ఇలా ఉంటుంది:
  \begin{prooftree}
    \AxiomC{$\Gamma$}
    \RightLabel{$\delta_1$}
    \DeduceC{$!A \land !B$}
    \RightLabel{\Elim{\land}}
    \UnaryInfC{$!A$}
  \end{prooftree}
  ఆగమన పరికల్పన ప్రకారం $\delta_1$లోని !!{undischarged} పరికల్పనలు
  $\Gamma$ నుంచి $!A \land !B$ అనుగమిస్తుంది. ఒక
  \iftag{FOL}{!!a{structure}~$\Struct{M}$ను}{సత్యనిర్దేశం~$\pAssign{v}$ను}
  పరిశీలించండి.
  \sourcecorrection{OLTEND-004}{మూలంలో “నిర్మాణం” అనే నామపదాన్ని
  ట్యాగ్ వెలుపల ఉంచడం వల్ల ప్రతిపాదనాత్మక రూపం “నిర్మాణం~$v$”గా
  వస్తుంది; FOL శాఖలో నిర్మాణం, మరొక శాఖలో సత్యనిర్దేశం వచ్చేలా
  సరిచేశాం.}
  $\iftag{FOL}{\Sat{M}}{\pSat{v}}{\Gamma}$ అయితే
  $\iftag{FOL}{\Sat{M}}{\pSat{v}}{!A}$ అని చూపాలి.
  $\iftag{FOL}{\Sat{M}}{\pSat{v}}{\Gamma}$ అనుకుందాం. మన ఆగమన
  పరికల్పన ($\Gamma \Entails !A \land !B$) ప్రకారం
  $\iftag{FOL}{\Sat{M}}{\pSat{v}}{!A \land !B}$. నిర్వచనం ప్రకారం
  $\iftag{FOL}{\Sat{M}}{\pSat{v}}{!A \land !B}$ అయితే, అప్పుడు మాత్రమే
  $\iftag{FOL}{\Sat{M}}{\pSat{v}}{!A}$ మరియు
  $\iftag{FOL}{\Sat{M}}{\pSat{v}}{!B}$. ($!A \land !B$ నుంచి $!B$ను
  నిగమించే సందర్భాన్ని కూడా ఇలాగే పరిష్కరిస్తాం.)

\item చివరి నిగమనం \Intro{\lor}: దీనికి రెండు రూపాలున్నాయి: పూర్వాధారం
  $!A$ నుంచి లేదా పూర్వాధారం $!B$ నుంచి $!A \lor !B$ను నిగమించవచ్చు.
  మొదటి సందర్భాన్ని పరిశీలిద్దాం. !!{derivation} రూపం ఇలా ఉంటుంది:
  \begin{prooftree}
    \AxiomC{$\Gamma$}
    \RightLabel{$\delta_1$}
    \DeduceC{$!A$}
    \RightLabel{\Intro{\lor}}
    \UnaryInfC{$!A \lor !B$}
  \end{prooftree}
  ఆగమన పరికల్పన ప్రకారం $\delta_1$లోని !!{undischarged} పరికల్పనలు
  $\Gamma$ నుంచి $!A$ అనుగమిస్తుంది. ఒక
  \iftag{FOL}{!!a{structure}~$\Struct{M}$}{!!a{valuation}~$\pAssign{v}$}ను
  పరిశీలించండి. $\iftag{FOL}{\Sat{M}}{\pSat{v}}{\Gamma}$ అయితే
  $\iftag{FOL}{\Sat{M}}{\pSat{v}}{!A \lor !B}$ అని చూపాలి.
  $\iftag{FOL}{\Sat{M}}{\pSat{v}}{\Gamma}$ అనుకుందాం; ఆగమన పరికల్పన
  $\Gamma \Entails !A$ కనుక $\iftag{FOL}{\Sat{M}}{\pSat{v}}{!A}$.
  అందువల్ల $\iftag{FOL}{\Sat{M}}{\pSat{v}}{!A \lor !B}$ కూడా.
  ($!B$ నుంచి $!A \lor !B$ను నిగమించే సందర్భాన్ని కూడా ఇలాగే
  పరిష్కరిస్తాం.)

\item చివరి నిగమనం \Intro{\lif}: పరికల్పన~$!A$, నిష్కర్ష~$!B$ గల
  ఉపనిరూపణ నుంచి $!A \lif !B$ను నిగమిస్తాం; అంటే
  \begin{prooftree}
    \AxiomC{$\Gamma, \Discharge{!A}{n}$}
    \RightLabel{$\delta_1$}
    \DeduceC{$!B$}
    \DischargeRule{\Intro{\lif}}{n}
    \UnaryInfC{$!A \lif !B$}
  \end{prooftree}
  ఆగమన పరికల్పన ప్రకారం $\delta_1$లోని !!{undischarged} పరికల్పనల
  నుంచి $!B$ అనుగమిస్తుంది; అంటే
  $\Gamma \cup \{!A\} \Entails !B$. ఒక
  \iftag{FOL}{!!a{structure}~$\Struct{M}$}{!!a{valuation}~$\pAssign{v}$}ను
  పరిశీలించండి. చివరి నిగమనం వద్ద $!A$ ఉపసంహరించబడింది కనుక
  $\delta$లోని !!{undischarged} పరికల్పనలు $\Gamma$ మాత్రమే. కాబట్టి
  $\Gamma \Entails !A \lif !B$ అని చూపాలి. ప్రతిపక్ష నిరూపణ కోసం,
  ఏదో \iftag{FOL}{!!{structure}~$\Struct{M}$}{!!{valuation}~$\pAssign{v}$}కు
  $\iftag{FOL}{\Sat{M}}{\pSat{v}}{\Gamma}$ కానీ
  $\iftag{FOL}{\Sat/{M}}{\pSat/{v}}{!A \lif !B}$ అనుకుందాం. అప్పుడు
  $\iftag{FOL}{\Sat{M}}{\pSat{v}}{!A}$ మరియు
  $\iftag{FOL}{\Sat/{M}}{\pSat/{v}}{!B}$. కానీ పరికల్పన ప్రకారం
  $!B$ అనేది $\Gamma \cup \{!A\}$ పర్యవసానం; అంటే
  $\iftag{FOL}{\Sat{M}}{\pSat{v}}{!B}$, ఇది వైరుధ్యం. కనుక
  $\Gamma \Entails !A \lif !B$.

\item చివరి నిగమనం \FalseInt: ఇక్కడ $\delta$ ఇలా ముగుస్తుంది:
  \begin{prooftree}
    \AxiomC{$\Gamma$}
    \RightLabel{$\delta_1$}
    \DeduceC{$\lfalse$}
    \RightLabel{\FalseInt}
    \UnaryInfC{$!A$}
  \end{prooftree}
  ఆగమన పరికల్పన ప్రకారం $\Gamma \Entails \lfalse$. మనం
  $\Gamma \Entails !A$ అని చూపాలి. అది కాదనుకుందాం; అప్పుడు ఏదో
  $\iftag{FOL}{\Struct{M}}{\pAssign{v}}$కు
    $\iftag{FOL}{\Sat{M}}{\pSat{v}}{\Gamma}$ మరియు
    $\iftag{FOL}{\Sat/{M}}{\pSat/{v}}{!A}$. కానీ ఎల్లప్పుడూ
    $\iftag{FOL}{\Sat/{M}}{\pSat/{v}}{\lfalse}$; కనుక ఇది
    $\Gamma \Entails/ \lfalse$ అవుతుందని అర్థం. అది ఆగమన పరికల్పనకు
    విరుద్ధం.

\item చివరి నిగమనం \FalseCl: అభ్యాసం.

\iftag{FOL}{
\item చివరి నిగమనం \Intro{\lforall}: అప్పుడు $\delta$ రూపం
  ఇలా ఉంటుంది:
  \begin{prooftree}
    \AxiomC{$\Gamma$}
    \RightLabel{$\delta_1$}
    \DeduceC{$!A(a)$}
    \RightLabel{\Intro{\lforall}}
    \UnaryInfC{$\lforall[x][!A(x)]$}
  \end{prooftree}
  ఆగమన పరికల్పన ప్రకారం పూర్వాధారం $!A(a)$ అనేది !!{undischarged}
  పరికల్పనలు $\Gamma$ యొక్క పర్యవసానం. $\Sat{M}{\Gamma}$ అయ్యే ఒక
  నిర్మాణం $\Struct{M}$ను పరిశీలించండి. మనం
  $\Sat{M}{\lforall[x][!A(x)]}$ అని చూపాలి. $\lforall[x][!A(x)]$
  ఒక !!a{sentence} కనుక ప్రతి చరరాశి నిర్దేశం~$s$కు
  $\Sat{M}{!A(x)}[s]$ అని చూపాలి
  (\olref[syn][ass]{prop:sat-quant}). $\Gamma$ అంతా వాక్యాలతోనే
  కూడి ఉంది కనుక ప్రతి $!B \in \Gamma$కు
  $\Sat{M}{!B}[s]$ (\olref[syn][sat]{defn:satisfaction}).
  $\Assign{a}{M'} = s(x)$ కావడం తప్ప $\Struct{M}$లాగే ఉండే
  $\Struct{M'}$ను తీసుకోండి. $a$ అనేది $\Gamma$లో కనిపించదు కనుక
  \olref[syn][ext]{cor:extensionality-sent} ప్రకారం
  $\Sat{M'}{\Gamma}$. $\Gamma \Entails !A(a)$ కనుక
  $\Sat{M'}{!A(a)}$. $!A(a)$ ఒక !!a{sentence} కనుక
  \olref[syn][ass]{prop:sentence-sat-true} ప్రకారం
  $\Sat{M'}{!A(a)}[s]$. $!A(a)$ అనేది
  $\Subst{!A(x)}{a}{x}$ మాత్రమేనని గుర్తుచేసుకుంటే,
  \olref[syn][ext]{prop:ext-formulas} ప్రకారం
  $\Sat{M'}{!A(x)}[s]$ అయితే, అప్పుడు మాత్రమే $\Sat{M'}{!A(a)}$.
  కనుక $\Sat{M'}{!A(x)}[s]$. $!A(x)$లో $a$ కనిపించదు కనుక
  \olref[syn][ext]{prop:extensionality} ప్రకారం
  $\Sat{M}{!A(x)}[s]$. అయితే $s$ను యాదృచ్ఛికంగా ఎంచుకున్నాం కనుక
  $\Sat{M}{\lforall[x][!A(x)]}$.

\item చివరి నిగమనం \Intro{\lexists}: అభ్యాసం.

\item చివరి నిగమనం \Elim{\forall}: అభ్యాసం.
}{}
\end{enumerate}

ఇప్పుడు అనేక పూర్వాధారాలు గల సాధ్యమైన నిగమనాలను పరిశీలిద్దాం:
\Elim{\lor}, \Intro{\land}, \iftag{FOL}{\Elim{\lif},
  \Elim{\lexists}}{\Elim{\lif}}.
\begin{enumerate}

\item చివరి నిగమనం \Intro{\land}. $!A$, $!B$ పూర్వాధారాల నుంచి
  $!A \land !B$ను నిగమిస్తాం; $\delta$ రూపం ఇలా ఉంటుంది:
  \begin{prooftree}
    \AxiomC{$\Gamma_1$}
    \RightLabel{$\delta_1$}
    \DeduceC{$!A$}
    \AxiomC{$\Gamma_2$}
    \RightLabel{$\delta_2$}
    \DeduceC{$!B$}
    \RightLabel{\Intro{\land}}
    \BinaryInfC{$!A \land !B$}
  \end{prooftree}
  ఆగమన పరికల్పన ప్రకారం $\delta_1$లోని !!{undischarged} పరికల్పనలు
  $\Gamma_1$ నుంచి $!A$ అనుగమిస్తుంది; $\delta_2$లోని
  !!{undischarged} పరికల్పనలు $\Gamma_2$ నుంచి $!B$ అనుగమిస్తుంది.
  $\delta$లోని !!{undischarged} పరికల్పనలు $\Gamma_1 \cup \Gamma_2$;
  కాబట్టి $\Gamma_1 \cup \Gamma_2 \Entails !A \land !B$ అని చూపాలి.
  $\iftag{FOL}{\Sat{M}}{\pSat{v}}{\Gamma_1 \cup \Gamma_2}$ అయ్యే
  ఒక \iftag{FOL}{!!a{structure}~$\Struct{M}$}{!!a{valuation}~$\pAssign{v}$}ను
  పరిశీలించండి. $\iftag{FOL}{\Sat{M}}{\pSat{v}}{\Gamma_1}$, మరియు
  $\Gamma_1 \Entails !A$ కనుక
  $\iftag{FOL}{\Sat{M}}{\pSat{v}}{!A}$. అలాగే
  $\iftag{FOL}{\Sat{M}}{\pSat{v}}{\Gamma_2}$, మరియు
  $\Gamma_2 \Entails !B$ కనుక $\iftag{FOL}{\Sat{M}}{\pSat{v}}{!B}$.
  రెండూ కలిపితే $\iftag{FOL}{\Sat{M}}{\pSat{v}}{!A \land !B}$.

\item చివరి నిగమనం \Elim{\lor}: అభ్యాసం.

\item చివరి నిగమనం \Elim{\lif}. $!A \lif !B$, $!A$ అనే
  పూర్వాధారాల నుంచి $!B$ను నిగమిస్తాం. !!{derivation}~$\delta$ ఇలా
  ఉంటుంది:
  \begin{prooftree}
    \AxiomC{$\Gamma_1$}
    \RightLabel{$\delta_1$}
    \DeduceC{$!A \lif !B$}
    \AxiomC{$\Gamma_2$}
    \RightLabel{$\delta_2$}
    \DeduceC{$!A$}
    \RightLabel{\Elim{\lif}}
    \BinaryInfC{$!B$}
  \end{prooftree}
  ఆగమన పరికల్పన ప్రకారం $\delta_1$లోని !!{undischarged} పరికల్పనలు
  $\Gamma_1$ నుంచి $!A \lif !B$ అనుగమిస్తుంది; $\delta_2$లోని
  !!{undischarged} పరికల్పనలు $\Gamma_2$ నుంచి $!A$ అనుగమిస్తుంది.
  ఒక \iftag{FOL}{!!a{structure}~$\Struct{M}$}{!!a{valuation}~$\pAssign{v}$}ను
  పరిశీలించండి. $\iftag{FOL}{\Sat{M}}{\pSat{v}}{\Gamma_1 \cup
    \Gamma_2}$ అయితే $\iftag{FOL}{\Sat{M}}{\pSat{v}}{!B}$ అని చూపాలి.
  $\iftag{FOL}{\Sat{M}}{\pSat{v}}{\Gamma_1 \cup \Gamma_2}$ అనుకుందాం.
  $\Gamma_1 \Entails !A \lif !B$ కనుక
  $\iftag{FOL}{\Sat{M}}{\pSat{v}}{!A \lif !B}$. అలాగే
  $\Gamma_2 \Entails !A$ కనుక $\iftag{FOL}{\Sat{M}}{\pSat{v}}{!A}$.
  అంటే $\iftag{FOL}{\Sat{M}}{\pSat{v}}{!B}$ (ఎందుకంటే
  $\iftag{FOL}{\Sat/{M}}{\pSat/{v}}{!B}$ అయితే,
  $\iftag{FOL}{\Sat{M}}{\pSat{v}}{!A}$ కనుక
  $\iftag{FOL}{\Sat/{M}}{\pSat/{v}}{!A \lif !B}$ అవుతుంది; అది
  $\iftag{FOL}{\Sat{M}}{\pSat{v}}{!A \lif !B}$కు విరుద్ధం).

\item చివరి నిగమనం \Elim{\lnot}: అభ్యాసం.

\tagitem{FOL}{చివరి నిగమనం \Elim{\lexists}: అభ్యాసం.}{}
\end{enumerate}
\end{proof}

\tagprob{FOL}
\begin{prob}
\olref[fol][ntd][sou]{thm:soundness} నిరూపణను పూర్తి చేయండి.
\end{prob}
\tagendprob

\tagprob{notFOL}
\begin{prob}
\olref[pl][ntd][sou]{thm:soundness} నిరూపణను పూర్తి చేయండి.
\end{prob}
\tagendprob

\begin{cor}
\ollabel{cor:weak-soundness}
$\Proves !A$ అయితే $!A$ ఒక \iftag{FOL}{చెల్లుబాటు అయ్యేది}{సర్వసత్యం}.
\end{cor}

\begin{cor}
\ollabel{cor:consistency-soundness}
$\Gamma$ సంతృప్తమైతే అది అవైరుధ్యమైనది.
\end{cor}

\begin{proof}
ప్రతిపక్షాన్ని నిరూపిస్తాం. $\Gamma$ అవైరుధ్యమైనది కాదనుకుందాం.
అప్పుడు $\Gamma \Proves \lfalse$; అంటే $\Gamma$లోని
!!{undischarged} పరికల్పనల నుంచి $\lfalse$ యొక్క ఒక !!a{derivation}
ఉంటుంది. \olref{thm:soundness} ప్రకారం $\Gamma$ను సంతృప్తిపరచే ఏ
\iftag{FOL}{!!{structure}~$\Struct{M}$}{!!{valuation}~$\pAssign{v}$}
అయినా $\lfalse$ను సంతృప్తిపరచాలి. ప్రతి
\iftag{FOL}{!!{structure}~$\Struct{M}$}{!!{valuation}~$\pAssign{v}$}కు
$\iftag{FOL}{\Sat/{M}}{\pSat/{v}}{\lfalse}$ కనుక ఏ
\iftag{FOL}{$\Struct{M}$}{$\pAssign{v}$} కూడా $\Gamma$ను
సంతృప్తిపరచలేదు; అంటే $\Gamma$ సంతృప్తం కాదు.
\end{proof}

\end{document}
